\scanpage{68}
% archive foliation: 67
\deleted{
$(p,q,r)$ doit être parmi les \deleted{\unsure{trois}} syst. énumérés

a) $p,2,2$ : $\dfrac1p+\dfrac12+\dfrac12>\left(1+\left(1-\dfrac1s\right)\right)$\quad i.e. $\dfrac1p+\dfrac1s>1$ impossible car $\dfrac1p,\dfrac1s\leq\dfrac12$

b) $(p,3,2)$\quad $(p=3,4,5)$, $s\geq p$ \deleted{\unsure{\ensuremath{\frac1p+\frac13+\frac12}}}
\[
\frac1s+\frac1p+\underbrace{\frac13+\frac12}_{5/6}>2\qquad
\frac1s+\frac1p>\frac76\qquad\text{impossible car }\frac1s,\frac1p\leq\frac13
\]
}
\[
\frac1p+\frac1q>2-\underbrace{\left(\frac1r+\frac1s\right)}_{\substack{\leq1\\\text{car }\frac1r\leq\frac12\\\frac1s\leq\frac12}}\geq1
\quad\text{donc}\quad\frac1p+\frac1q\geq1\quad\text{impossible car }\frac1p,\frac1q\leq\frac12
\]
\uline{Impossible} !

Car si $\operatorname{Card}I\geq4$, on aura
\[
\frac1{d_1}+\frac1{d_2}+\frac1{d_3}+\frac1{d_4}
\deleted{\unsure{+\sum_{5\leq i\leq\nu}\frac1{d_i}}}
>2+\sum_{5\leq i\leq\nu}\left(1-\frac1{d_i}\right)
\underset{2}{\unsure{\vee|\geq}}
\]
et on vient de voir que c'est impossible !

\deleted{
\uline{Étude du cas} \deleted{$\operatorname{Card}I=2$}\inserted{$(N,N)$} : En caractéristique $0$, c'est le cas bien connu
\[
G\simeq\mu_N(k)\simeq\mathbb Z/N\mathbb Z
\]
\deleted{\unsure{Plus gén\textsuperscript{t}|Plus généralement}} \unsure{gén.|Bien}, il en est ainsi si car. $p$ premier à $N$.

Si $p\mid N$, on trouve une $p$-extension \inserted{(de degré $p^r$, tlt ramifiée et \unsure{inséparable})} d'une $N'$-extension kummérienne $(N=p^rN',\ (p,N')=1)$. C'est donc une \uline{extension} \deleted{et} \uline{résoluble}, \inserted{mais mieux :} \deleted{\unsure{On ne peut être plus précis},} \deleted{de} \inserted{les $p$-\deleted{extensions} \unsure{sous-}}

\inserted{\unsure{groupes} $H$ de $\operatorname{GP}(1,k)$ sont \uline{abéliens} et contenus dans}
\deleted{\unsure{toute hypothèse|toute la théorie}. \unsure{On déduit même que l'extension est}}
\deleted{\uline{abélienne finie}}
sous-groupe unipotent de dim $1$, \inserted{conjugués entre eux, et $H$ isom : $(\mathbb Z/p\mathbb Z)^{\unsure{\nu|r}}$.} On trouve donc

$G\subset\text{Borel}$, et par suite $G$ produit $\frac12$ direct d'un groupe cyclique premier à \unsure{$2$|$p$} car., par un $(\mathbb Z/p\mathbb Z)^r$. En fait, on voit aisément (\deleted{\unsure{ici encore}} qu'on est ds le cas précédent (cyclique premier à car., donc kummérien, qui en est exclu), on \unsure{tombe} en contradiction.

\margin{En tous les cas, on a une extension abélienne, \unsure{non kummérienne|kummérienne} (d'ordre premier à $p$) \unsure{sous} \deleted{\unsure{gr.}} \unsure{gr.} $(\mathbb Z/p^r\mathbb Z)^\nu$}

\uline{Étude du cas} $(N,2,2)$, $(\#G=2N)$. En car. $\ne2$, \uline{l'indice}\inserted{\unsure{d'un pt}} de ramification $2$ est ramification modérée, et faisant une \deleted{\unsure{différence}} extension kummérienne de degré $2$, on \deleted{\unsure{tr}} trouve un $X'/Y'$ qui n'a plus que $2$ pts de ramification \deleted{\unsure{dans}} $Y'$, donc $X'/Y'$ est justiciable de l'étude précédente mais \deleted{\unsure{il faut justifier que|il faut vérifier que}}\inserted{\unsure{voir pour|voir}} si $X'$ de genre $0$ (facile)
\diagram{\begin{tikzcd}[ampersand replacement=\&]
X \arrow[d] \& X' \arrow[l] \arrow[d,dash]\\
Y \& Y' \arrow[l]
\end{tikzcd}}
}
